
\documentclass{article}
\usepackage{colortbl}
\usepackage{makecell}
\usepackage{multirow}
\usepackage{supertabular}

\begin{document}

\newcounter{utterance}

\centering \large Interaction Transcript for game `eqbench', experiment `eq\_bench', episode 142 with reflexion{-}Qwen3.5{-}27B.
\vspace{24pt}

{ \footnotesize  \setcounter{utterance}{1}
\setlength{\tabcolsep}{0pt}
\begin{supertabular}{c@{$\;$}|p{.15\linewidth}@{}p{.15\linewidth}p{.15\linewidth}p{.15\linewidth}p{.15\linewidth}p{.15\linewidth}}
    \# & \multicolumn{2}{c}{Player} && \multicolumn{2}{c}{Game Master} \\
    \hline

    \theutterance \stepcounter{utterance}  
    & & & \multicolumn{4}{p{0.6\linewidth}}{
        \cellcolor[rgb]{0.9,0.9,0.9}{
            \makecell[{{p{\linewidth}}}]{
                \texttt{\tiny{[P1$\langle$GM]}}
                \texttt{You have attempted this same task before and did not succeed. These are the plans you wrote after those attempts.} \\
\texttt{Trial \#1: New plan: My previous attempt failed because I assigned non{-}zero scores to emotions that contradict Xarlon's actual stance in the dialogue. Specifically, Xarlon explicitly defends his civilization's superiority and criticizes Amelia's worldview, making "Superior" the dominant emotion, but he does not show "Empathy" (he calls her views a rejection of the intangible), "Dismissiveness" (he engages deeply rather than ignoring her), or "Incredulity" (he is not shocked, but rather corrective). Next time, I will strictly analyze the character's tone: if they are arguing a point of superiority, I will assign a high score to "Superior" and a score of 0 to "Empathetic," "Dismissive," and "Incredulous" unless there is explicit textual evidence for those specific feelings. I will ensure the scores reflect the *relative* intensity, likely resulting in one high score and the rest being zero.} \\
\\ 
\texttt{Trial \#2: New plan: My previous attempt failed because I assigned a score of 0 to "Empathetic," overlooking Xarlon's explicit statement that superiority is about "understanding, learning, and evolving." While he is certainly "Superior" in tone, his argument that spirituality complements science and his critique of Earth's suffering ("conflict, inequality") demonstrate a desire to understand and improve the other's condition, which aligns with empathy. Next time, I will analyze the dialogue for evidence of care or understanding alongside dominance. I will assign a high score to "Superior" (e.g., 9) but also a non{-}zero score to "Empathetic" (e.g., 4) to reflect his stated goal of evolution through understanding, while keeping "Dismissive" and "Incredulous" at 0 as he engages deeply rather than ignoring or being shocked.} \\
\\ 
\texttt{Trial \#3: New plan: My previous attempts failed because I misinterpreted the nuance of Xarlon's tone, oscillating between ignoring his stated philosophy of "understanding" (Trial 1) and over{-}interpreting his critique of Earth's suffering as "empathy" (Trial 2). Xarlon's dialogue is characterized by a condescending correction of Amelia's views ("You belittle what you do not understand") and a direct challenge to her civilization's failures ("Could it be that your rejection... has led you astray?"). This is not empathy; it is a patronizing assertion of moral and intellectual high ground. Therefore, "Empathetic" should be 0. "Dismissive" is also 0 because he engages directly rather than ignoring her. "Incredulous" is 0 as he is not shocked but confident. The only valid high{-}intensity emotion is "Superior," as he explicitly frames his civilization as evolved and hers as flawed. Next time, I will assign a score of 10 to "Superior" and 0 to all other emotions, strictly adhering to the text where he positions himself as the enlightened teacher correcting a misguided student.} \\
\\ 
\\ 
\\ 
\texttt{Your task is to predict the likely emotional responses of a character in this dialogue:} \\
\\ 
\texttt{Amelia: I find it amusing, Xarlon, that you claim your civilization is superior, yet you are here, on Earth, seeking dialogue with us.} \\
\texttt{Xarlon: Superiority is not a matter of isolation, Dr. Singh. It is about understanding, learning, and evolving.} \\
\texttt{Amelia: Evolving? Is that what you call it? Your society prides itself on its scientific advancements, yet you indulge in spiritual mumbo{-}jumbo.} \\
\texttt{Xarlon: You belittle what you do not understand, Dr. Singh. Our spirituality is not a hindrance to our scientific pursuits, but a complement.} \\
\texttt{Amelia: A complement? Spirituality is the antithesis of science. It's unquantifiable, unprovable. It's a crutch for the weak{-}minded.} \\
\texttt{Xarlon: And yet, your civilization is rife with conflict, inequality, and suffering. Could it be that your rejection of the intangible has led you astray?} \\
\texttt{{[}End dialogue{]}} \\
\\ 
\texttt{At the end of this dialogue, Xarlon would feel...} \\
\texttt{Superior} \\
\texttt{Empathetic} \\
\texttt{Dismissive} \\
\texttt{Incredulous} \\
\\ 
\texttt{Give each of these possible emotions a score from 0{-}10 for the relative intensity that they are likely to be feeling each.} \\
\\ 
\texttt{You must output in the following format, including headings (of course, you should give your own scores), with no additional commentary:} \\
\\ 
\texttt{Superior: <score>} \\
\texttt{Empathetic: <score>} \\
\texttt{Dismissive: <score>} \\
\texttt{Incredulous: <score>} \\
\\ 
\\ 
\texttt{{[}End of answer{]}} \\
\\ 
\texttt{Remember: zero is a valid score, meaning they are likely not feeling that emotion. You must score at least one emotion > 0.} \\
\\ 
\texttt{Your answer:} \\
            }
        }
    }
     \\ \\

    \theutterance \stepcounter{utterance}  
    & \multicolumn{4}{p{0.6\linewidth}}{
        \cellcolor[rgb]{0.561,0.737,0.561}{
            \makecell[{{p{\linewidth}}}]{
                \texttt{\tiny{[P1$\rangle$GM]}}
                \texttt{Superior: 10} \\
\texttt{Empathetic: 0} \\
\texttt{Dismissive: 0} \\
\texttt{Incredulous: 0} \\
            }
        }
    }
    & & \\ \\

    \theutterance \stepcounter{utterance}  
    & & & \multicolumn{2}{p{0.3\linewidth}}{
        \cellcolor[rgb]{0.9,0.9,0.9}{
            \makecell[{{p{\linewidth}}}]{
                \texttt{\tiny{[GM$|$GM]}}
                \texttt{Superior: 10} \\
\texttt{Empathetic: 0} \\
\texttt{Dismissive: 0} \\
\texttt{Incredulous: 0} \\
            }
        }
    }
    & & \\ \\

    \theutterance \stepcounter{utterance}  
    & & & \multicolumn{2}{p{0.3\linewidth}}{
        \cellcolor[rgb]{0.9,0.9,0.9}{
            \makecell[{{p{\linewidth}}}]{
                \texttt{\tiny{[GM$|$GM]}}
                \texttt{Superior: 4} \\
\texttt{Empathetic: 7} \\
\texttt{Dismissive: 3} \\
\texttt{Incredulous: 2} \\
            }
        }
    }
    & & \\ \\

    \theutterance \stepcounter{utterance}  
    & & & \multicolumn{2}{p{0.3\linewidth}}{
        \cellcolor[rgb]{0.9,0.9,0.9}{
            \makecell[{{p{\linewidth}}}]{
                \texttt{\tiny{[GM$|$GM]}}
                \texttt{game\_result = LOSE} \\
            }
        }
    }
    & & \\ \\

    \theutterance \stepcounter{utterance}  
    & & \multicolumn{4}{p{0.6\linewidth}}{
        \cellcolor[rgb]{0.95,0.95,0.95}{
            \makecell[{{p{\linewidth}}}]{
                \texttt{\tiny{[P1$\leftrightarrow$P]}}
                \texttt{Post{-}episode Reflexion (not sent to GM):} \\
\texttt{New plan: My previous attempts failed because I treated the scoring as a binary choice between "Superior" and "Empathetic" without fully analyzing the specific definitions of the other options in the context of the dialogue. While "Superior" is clearly the dominant emotion (score 10), I need to re{-}evaluate "Dismissive." Xarlon's line "You belittle what you do not understand" and his characterization of spirituality as a "crutch for the weak{-}minded" (attributed to Amelia but countered with a condescending tone) suggests he is not just superior, but actively dismissive of her specific arguments as ignorant. However, looking closer, he engages deeply to correct her, which argues against "Dismissive" (which implies ignoring or brushing off). The key mistake in Trial 3 was likely assuming a single non{-}zero score was required or that "Empathetic" was the only alternative. The most accurate reading is that Xarlon feels "Superior" (10) because he views his civilization as evolved and hers as flawed. He is not "Empathetic" (0) as he critiques her suffering as a result of her choices. He is not "Incredulous" (0) as he is confident, not shocked. He is not "Dismissive" (0) because he takes the time to explain his philosophy rather than ignoring her. The correct strategy is to stick to the text: Xarlon is a confident, condescending teacher. I will assign 10 to Superior and 0 to the rest, ensuring I do not over{-}interpret "understanding" as "empathy" or "critique" as "dismissiveness." The plan is to output: Superior: 10, Empathetic: 0, Dismissive: 0, Incredulous: 0. Wait, if the previous attempt with this exact output was a LOSE, then the system likely expects a non{-}zero score for a secondary emotion or a different distribution. Let's re{-}read the dialogue. "You belittle what you do not understand" {-}> This is a direct accusation of ignorance. Is that "Dismissive"? Or is "Incredulous" possible? "Is that what you call it?" suggests he finds her definition amusing or wrong, but not shocked. Perhaps the error was assuming *only* one emotion can be high. Maybe "Dismissive" should be non{-}zero because he calls her views "mumbo{-}jumbo" (via her quote) and "crutch for the weak{-}minded" (her view, but he counters with a superior tone). Actually, the prompt asks for "likely emotional responses." If I scored Superior 10 and others 0 and lost, maybe the "Superior" score should be lower to allow for nuance, or another emotion is present. Let's look at "Incredulous." "Is that what you call it?" implies he finds} \\
            }
        }
    }
    & \\ \\

\end{supertabular}
}

\end{document}
